% ============================================================================
% 03_optimization/01_algorithm.tex
% Section body for inclusion immediately after the corresponding \section
% command in the master book file. It follows section_body_template.tex.
% No document-level commands, package imports, \chapter, or \section commands.
% ============================================================================

The preceding discussion established the mathematical components from which a
language model is constructed. Those components define a parameterized
computational function, but they do not by themselves determine the numerical
values of its trainable parameters. Training requires a procedure that uses
observed data to modify those parameters so that the model's predictions
become progressively more consistent with the training objective. This
procedure is the subject of optimization. In a large language model (LLM),
optimization is not merely a numerical routine appended to the neural
architecture: it is a central part of the training system, because its update
rule determines how information contained in the training data is translated
into changes in the model.

The central problem is to minimize a loss function over a parameter space
whose dimension may reach billions. The objective is generally nonconvex,
gradients are usually estimated from minibatches rather than the complete
training distribution, and the optimization process must operate within
constraints imposed by accelerator memory, numerical precision, computation,
and communication. These properties explain why the development of LLM
optimizers involves more than selecting a learning rate. It requires decisions
about which gradient information to retain, how to transform that information
into parameter updates, what approximation to the local geometry is affordable,
and how to evaluate progress under a fixed token or wall-clock budget
\parencite{Bottou2018Optimization,Ranganath2026Navigating}.

\subsection{Optimization as a Statistical Learning Problem}
\label{subsec:optimization-statistical-learning}

Let $\boldsymbol{\theta}\in\mathbb{R}^{d}$ denote the complete vector of
trainable parameters of a language model, where $d$ is the number of scalar
parameters. For an autoregressive language model and a token sequence
$x_{1:T}=(x_1,\ldots,x_T)$, the joint probability is factorized into
conditional next-token probabilities:

\begin{equation}
    p_{\boldsymbol{\theta}}(x_{1:T})
    =
    \prod_{t=1}^{T}
    p_{\boldsymbol{\theta}}(x_t\mid x_{<t}),
    \label{eq:opt-autoregressive-factorization}
\end{equation}

where $x_{<t}=(x_1,\ldots,x_{t-1})$ is the context preceding token $x_t$.
Training by next-token prediction minimizes the expected negative
log-likelihood under the training distribution $\mathcal{D}$:

\begin{equation}
    F(\boldsymbol{\theta})
    =
    -\mathbb{E}_{x_{1:T}\sim\mathcal{D}}
    \left[
        \sum_{t=1}^{T}
        \log p_{\boldsymbol{\theta}}(x_t\mid x_{<t})
    \right].
    \label{eq:opt-expected-nll}
\end{equation}

This objective is an instance of the expected-risk formulation used in
statistical learning. More generally, let $\xi$ denote a random training
observation, which may be a single example or a minibatch, and let
$f(\boldsymbol{\theta};\xi)$ be the loss evaluated at that observation. The
expected objective is

\begin{equation}
    F(\boldsymbol{\theta})
    =
    \mathbb{E}_{\xi\sim\mathcal{D}}
    \left[
        f(\boldsymbol{\theta};\xi)
    \right].
    \label{eq:opt-expected-risk}
\end{equation}

The distribution $\mathcal{D}$ is not known in complete form. Instead, the
training procedure has access to a finite dataset
$\{\xi_1,\ldots,\xi_N\}$. The corresponding empirical objective is

\begin{equation}
    F_N(\boldsymbol{\theta})
    =
    \frac{1}{N}
    \sum_{i=1}^{N}
    f(\boldsymbol{\theta};\xi_i),
    \label{eq:opt-empirical-risk}
\end{equation}

and the finite-data optimization problem is

\begin{equation}
    \boldsymbol{\theta}_N^{*}
    \in
    \underset{\boldsymbol{\theta}\in\mathbb{R}^{d}}
    {\operatorname{argmin}}
    F_N(\boldsymbol{\theta}).
    \label{eq:opt-empirical-minimization}
\end{equation}

The expected risk and empirical risk represent different mathematical
objects. The former measures expected loss under the underlying distribution;
the latter is the average loss over the available observations. Minimizing
the empirical objective is therefore a practical approximation to the
statistical learning problem, not a guarantee that the expected risk has
also been minimized. The distinction between expected and empirical risk
is the starting point for the optimization analysis of large-scale machine
learning \parencite{Bottou2018Optimization}. The autoregressive formulation
in \cref{eq:opt-autoregressive-factorization,eq:opt-expected-nll} makes this
general framework specific to language-model pretraining
\parencite{Ranganath2026Navigating}.

An optimization algorithm constructs a sequence of parameter vectors
$\boldsymbol{\theta}_1,\boldsymbol{\theta}_2,\ldots$ with the intention of
reducing the objective. The most basic local information used by
gradient-based methods is the gradient:

\begin{equation}
    \nabla F(\boldsymbol{\theta})
    =
    \begin{bmatrix}
        \frac{\partial F}{\partial\theta_1} &
        \frac{\partial F}{\partial\theta_2} &
        \cdots &
        \frac{\partial F}{\partial\theta_d}
    \end{bmatrix}^{\mathsf{T}}.
    \label{eq:opt-gradient}
\end{equation}

The gradient belongs to the same $d$-dimensional vector space as the
parameters. Its $j$-th component measures the local sensitivity of the
objective to the $j$-th parameter. For a sufficiently small displacement
$\Delta\boldsymbol{\theta}$, the first-order Taylor approximation gives

\begin{equation}
    F(\boldsymbol{\theta}+\Delta\boldsymbol{\theta})
    \approx
    F(\boldsymbol{\theta})
    +
    \nabla F(\boldsymbol{\theta})^{\mathsf{T}}
    \Delta\boldsymbol{\theta}.
    \label{eq:opt-first-order-taylor}
\end{equation}

Choosing the displacement in the negative-gradient direction,

\begin{equation}
    \Delta\boldsymbol{\theta}
    =
    -\eta\nabla F(\boldsymbol{\theta}),
    \qquad \eta>0,
    \label{eq:opt-gradient-displacement}
\end{equation}

yields the first-order change

\begin{equation}
    F(\boldsymbol{\theta}+\Delta\boldsymbol{\theta})
    -
    F(\boldsymbol{\theta})
    \approx
    -\eta\|\nabla F(\boldsymbol{\theta})\|_2^2.
    \label{eq:opt-first-order-decrease}
\end{equation}

Thus, the negative gradient is a local descent direction whenever the
gradient is nonzero. The scalar $\eta$ is the step size or learning rate;
it controls the magnitude of the displacement. This local argument
motivates gradient descent, but it does not by itself establish convergence.
A convergence analysis requires assumptions about how rapidly the gradient
can change.

A common assumption is that the gradient is Lipschitz continuous. A
differentiable function has an $L$-Lipschitz gradient if

\begin{equation}
    \|\nabla F(\boldsymbol{\theta})
      -\nabla F(\boldsymbol{\theta}')\|_2
    \leq
    L\|\boldsymbol{\theta}-\boldsymbol{\theta}'\|_2
    \label{eq:opt-lipschitz-gradient}
\end{equation}

for all relevant parameter vectors. Under this assumption, the descent
lemma states that

\begin{equation}
    F(\boldsymbol{\theta}')
    \leq
    F(\boldsymbol{\theta})
    +
    \nabla F(\boldsymbol{\theta})^{\mathsf{T}}
    (\boldsymbol{\theta}'-\boldsymbol{\theta})
    +
    \frac{L}{2}
    \|\boldsymbol{\theta}'-\boldsymbol{\theta}\|_2^2.
    \label{eq:opt-descent-lemma}
\end{equation}

Substituting the gradient-descent update
$\boldsymbol{\theta}_{t+1}
=\boldsymbol{\theta}_t-\eta_t\nabla F(\boldsymbol{\theta}_t)$
into \cref{eq:opt-descent-lemma} gives

\begin{align}
    F(\boldsymbol{\theta}_{t+1})
    &\leq
    F(\boldsymbol{\theta}_t)
    -
    \eta_t\|\nabla F(\boldsymbol{\theta}_t)\|_2^2
    +
    \frac{L\eta_t^2}{2}
    \|\nabla F(\boldsymbol{\theta}_t)\|_2^2
    \nonumber\\
    &=
    F(\boldsymbol{\theta}_t)
    -
    \eta_t\left(1-\frac{L\eta_t}{2}\right)
    \|\nabla F(\boldsymbol{\theta}_t)\|_2^2.
    \label{eq:opt-gradient-descent-bound}
\end{align}

For $0<\eta_t\leq 1/L$, the coefficient multiplying the squared gradient
is positive, and the bound guarantees a decrease in the objective whenever
the gradient is nonzero. This explains why an excessively large learning
rate can destabilize optimization: the quadratic remainder in the descent
bound can offset the improvement predicted by the linear term. An
excessively small learning rate, on the other hand, can make progress
unnecessarily slow \parencite{Bottou2018Optimization}.

The optimization landscape of a deep neural network is generally nonconvex.
Consequently, a stationary point satisfying
$\nabla F(\boldsymbol{\theta})=\mathbf{0}$ need not be a global minimizer.
It may be a local minimum, a saddle point, or another stationary point.
Theoretical results for nonconvex stochastic optimization therefore often
establish convergence toward regions with small gradient norm rather than
convergence to a global optimum. This distinction is essential when
interpreting guarantees for the training of Transformer language models
\parencite{Bottou2018Optimization,Ranganath2026Navigating}.

\subsection{Stochastic and Batch Gradient Methods}
\label{subsec:optimization-stochastic-batch}

Differentiating the empirical objective in
\cref{eq:opt-empirical-risk} gives the full gradient

\begin{equation}
    \nabla F_N(\boldsymbol{\theta})
    =
    \frac{1}{N}
    \sum_{i=1}^{N}
    \nabla f(\boldsymbol{\theta};\xi_i).
    \label{eq:opt-full-gradient}
\end{equation}

A batch-gradient method evaluates this complete sum at every iteration:

\begin{equation}
    \boldsymbol{\theta}_{t+1}
    =
    \boldsymbol{\theta}_t
    -
    \eta_t\nabla F_N(\boldsymbol{\theta}_t).
    \label{eq:opt-batch-gradient}
\end{equation}

The direction is calculated from all training observations and is therefore
the exact gradient of the empirical objective. Its disadvantage is
computational: each update requires processing the complete dataset. When
the dataset contains an enormous number of sequences and tokens, this
requirement makes frequent full-gradient evaluations impractical.

Stochastic gradient (SG) methods replace the full gradient with an estimate
computed from a randomly selected observation. If $i_t$ is sampled uniformly
from $\{1,\ldots,N\}$, the stochastic gradient is

\begin{equation}
    \mathbf{g}_t
    =
    \nabla f(\boldsymbol{\theta}_t;\xi_{i_t}),
    \label{eq:opt-stochastic-gradient}
\end{equation}

and the update is

\begin{equation}
    \boxed{
    \boldsymbol{\theta}_{t+1}
    =
    \boldsymbol{\theta}_t-\eta_t\mathbf{g}_t.
    }
    \label{eq:opt-sgd-update}
\end{equation}

Under uniform sampling, the estimator is conditionally unbiased:

\begin{equation}
    \mathbb{E}
    [\mathbf{g}_t\mid\boldsymbol{\theta}_t]
    =
    \nabla F_N(\boldsymbol{\theta}_t).
    \label{eq:opt-unbiased-gradient}
\end{equation}

Unbiasedness does not mean that every sampled direction decreases the
empirical objective. An individual stochastic gradient may point in a
direction that increases the full objective. It means that the expected
direction, conditioned on the current parameters, equals the full
empirical gradient. The iterates are therefore a stochastic process whose
trajectory depends on the sampled observations
\parencite{Bottou2018Optimization}.

The difference between the stochastic gradient and its conditional
expectation is gradient noise. Its covariance matrix is

\begin{equation}
    \boldsymbol{\Sigma}_t
    =
    \mathbb{E}\left[
       (\mathbf{g}_t-\nabla F_N)
       (\mathbf{g}_t-\nabla F_N)^{\mathsf{T}}
       \mid\boldsymbol{\theta}_t
    \right],
    \label{eq:opt-gradient-covariance}
\end{equation}

and its total variance is

\begin{equation}
    \sigma_t^2
    =
    \operatorname{Tr}(\boldsymbol{\Sigma}_t)
    =
    \mathbb{E}\left[
       \|\mathbf{g}_t-\nabla F_N\|_2^2
       \mid\boldsymbol{\theta}_t
    \right].
    \label{eq:opt-gradient-variance}
\end{equation}

The diagonal entries of $\boldsymbol{\Sigma}_t$ describe coordinate-wise
variances, while its off-diagonal entries describe correlations between
gradient coordinates. In high-dimensional neural networks, these
correlations need not vanish. The covariance therefore contains information
about the statistical structure of gradient noise, not merely the magnitude
of each individual coordinate.

For an unbiased estimator,

\begin{equation}
    \mathbb{E}
    [\|\mathbf{g}_t\|_2^2\mid\boldsymbol{\theta}_t]
    =
    \|\nabla F_N(\boldsymbol{\theta}_t)\|_2^2+\sigma_t^2.
    \label{eq:opt-second-moment-decomposition}
\end{equation}

Combining this identity with the descent lemma gives, under the same
smoothness assumptions,

\begin{align}
    \mathbb{E}
    [F_N(\boldsymbol{\theta}_{t+1})\mid\boldsymbol{\theta}_t]
    \leq\;&
    F_N(\boldsymbol{\theta}_t)
    -
    \eta_t\left(1-\frac{L\eta_t}{2}\right)
    \|\nabla F_N(\boldsymbol{\theta}_t)\|_2^2
    \nonumber\\
    &+
    \frac{L\eta_t^2}{2}\sigma_t^2.
    \label{eq:opt-stochastic-descent-bound}
\end{align}

The negative term represents expected progress in the gradient direction;
the positive term represents the adverse contribution of stochastic
variance. This inequality makes explicit why learning-rate selection and
noise control are central to stochastic optimization.

In large-scale neural-network training, gradients are commonly estimated
using a minibatch. For a minibatch $\mathcal{B}_t$ containing $B$ sampled
observations,

\begin{equation}
    \mathbf{g}_t
    =
    \frac{1}{B}
    \sum_{i\in\mathcal{B}_t}
    \nabla f(\boldsymbol{\theta}_t;\xi_i).
    \label{eq:opt-minibatch-gradient}
\end{equation}

When observations are sampled independently and have finite variance
$\sigma^2$, the variance of the minibatch mean is

\begin{equation}
    \operatorname{Var}(\mathbf{g}_t)
    =
    \frac{\sigma^2}{B}.
    \label{eq:opt-minibatch-variance}
\end{equation}

Increasing $B$ therefore reduces the variance of the gradient estimate.
However, it also increases the work performed to calculate one update.
The reduction in noise must be considered together with the computational
cost of processing a larger minibatch. Minibatches also provide an
important systems advantage: they expose parallel computation and allow
accelerators to process many training examples efficiently
\parencite{Bottou2018Optimization,Ranganath2026Navigating}.

For strongly convex objectives, stochastic-gradient methods with a constant
step size generally approach a neighborhood of the minimizer when gradient
noise remains nonzero. The neighborhood reflects a balance between
contraction toward the optimum and the variance introduced by stochastic
updates. A diminishing step-size sequence can reduce the asymptotic effect
of this noise. A classical sufficient condition is

\begin{equation}
    \sum_{t=1}^{\infty}\eta_t=\infty,
    \qquad
    \sum_{t=1}^{\infty}\eta_t^2<\infty.
    \label{eq:opt-robbins-monro-steps}
\end{equation}

The first condition prevents the cumulative step sizes from becoming
finite, while the second controls the accumulated effect of stochastic
noise. Under the assumptions developed by Bottou, Curtis, and Nocedal,
diminishing step sizes yield an $O(1/t)$ expected optimality-gap rate for
strongly convex objectives \parencite{Bottou2018Optimization}.

For nonconvex objectives, the relevant theoretical quantity is often the
squared gradient norm. Under the assumptions of the stochastic-gradient
analysis in \textcite{Bottou2018Optimization}, a fixed sufficiently small
step size gives a bound of the general form

\begin{equation}
    \frac{1}{K}
    \sum_{t=1}^{K}
    \mathbb{E}
    \left[
        \|\nabla F(\boldsymbol{\theta}_t)\|_2^2
    \right]
    \leq
    C_1\eta+
    \frac{C_2}{K\eta},
    \label{eq:opt-nonconvex-stationarity}
\end{equation}

where $C_1$ depends on the noise and smoothness bounds and $C_2$ depends
on the initial objective gap and the lower bound on the objective. The
first term reflects the residual effect of stochastic noise; the second
decreases with the number of iterations. This result does not guarantee a
global minimum. It shows that the algorithm spends an increasing fraction
of its iterations in regions where the objective has a relatively small
gradient.

\paragraph{Noise-reduction methods.}
The variance in a stochastic direction can be reduced without abandoning
the computational advantages of stochastic optimization. Dynamic
sample-size methods increase the number of observations used to estimate
the gradient as training progresses. They can use inexpensive estimates
during early optimization and more accurate estimates when the algorithm
approaches a stationary region.

Gradient-aggregation methods retain information from earlier iterations.
For example, the stochastic variance-reduced gradient (SVRG) estimator
uses a reference point $\widetilde{\boldsymbol{\theta}}$ and the full
gradient at that reference point:

\begin{equation}
    \mathbf{g}_t^{\mathrm{SVRG}}
    =
    \nabla f_{i_t}(\boldsymbol{\theta}_t)
    -
    \nabla f_{i_t}(\widetilde{\boldsymbol{\theta}})
    +
    \nabla F_N(\widetilde{\boldsymbol{\theta}}).
    \label{eq:opt-svrg}
\end{equation}

Taking the conditional expectation over $i_t$ gives
$\mathbb{E}[\mathbf{g}_t^{\mathrm{SVRG}}]
=\nabla F_N(\boldsymbol{\theta}_t)$.
The correction term removes the sampled gradient contribution evaluated
at the reference point and replaces it with the full reference gradient.
When the current iterate is close to the reference point, this construction
can substantially reduce variance. SAGA and stochastic average gradient
(SAG) are related gradient-aggregation approaches that retain historical
gradient information. Iterate averaging is another variance-reduction
strategy; it forms an average of iterates rather than relying only on the
most recent parameter vector:

\begin{equation}
    \overline{\boldsymbol{\theta}}_T
    =
    \frac{1}{T}
    \sum_{t=1}^{T}\boldsymbol{\theta}_t.
    \label{eq:opt-iterate-average}
\end{equation}

These methods are important in the broader theory of large-scale
optimization, although their memory and data-access requirements must be
considered before applying them to very large language-model training
\parencite{Bottou2018Optimization}.

\subsection{First-Order and Adaptive Optimization Algorithms}
\label{subsec:optimization-first-order}

A first-order optimizer constructs its update primarily from gradient
information rather than explicitly forming and inverting the Hessian.
A general update can be written as

\begin{equation}
    \boldsymbol{\theta}_{t+1}
    =
    \boldsymbol{\theta}_t-\eta_t\mathbf{u}_t,
    \label{eq:opt-general-update}
\end{equation}

where $\mathbf{u}_t$ is a direction derived from the current gradient and,
possibly, information retained from previous iterations. The algorithms
in this family differ in how they construct $\mathbf{u}_t$ and how much
additional state they store.

\paragraph{Stochastic gradient descent and momentum.}
The simplest first-order method uses the stochastic gradient directly:

\begin{equation}
    \boldsymbol{\theta}_{t+1}
    =
    \boldsymbol{\theta}_t-\eta_t\mathbf{g}_t.
    \label{eq:opt-sgd}
\end{equation}

Its principal advantage is its small memory overhead: it does not require
additional per-parameter optimizer state. Its limitation is that one
scalar learning rate is applied to all coordinates. In an LLM, embeddings,
attention projections, feed-forward matrices, and normalization parameters
can have different gradient statistics, so a single global learning rate
may not be equally appropriate for every parameter group
\parencite{Ranganath2026Navigating}.

Momentum introduces a velocity vector that accumulates gradient
information:

\begin{align}
    \mathbf{v}_t
    &=
    \beta\mathbf{v}_{t-1}+\mathbf{g}_t,
    \label{eq:opt-momentum-state}\\
    \boldsymbol{\theta}_{t+1}
    &=
    \boldsymbol{\theta}_t-\eta_t\mathbf{v}_t.
    \label{eq:opt-momentum-update}
\end{align}

Here, $\beta\in[0,1)$ controls the influence of previous gradients. By
smoothing short-term fluctuations, momentum can reduce oscillation and
improve progress in poorly conditioned landscapes. It does not, however,
provide coordinate-wise learning-rate adaptation
\parencite{Bottou2018Optimization,Ranganath2026Navigating}.

\paragraph{Adaptive diagonal preconditioning.}
Adaptive optimizers can be interpreted as diagonal preconditioned
gradient methods. A diagonal preconditioner rescales each coordinate
independently, without representing cross-coordinate correlations.

AdaGrad accumulates squared gradients:

\begin{align}
    \mathbf{s}_t
    &=
    \mathbf{s}_{t-1}
    +
    \mathbf{g}_t\odot\mathbf{g}_t,
    \label{eq:opt-adagrad-state}\\
    \boldsymbol{\theta}_{t+1}
    &=
    \boldsymbol{\theta}_t
    -
    \eta_t
    \frac{\mathbf{g}_t}{\sqrt{\mathbf{s}_t}+\epsilon}.
    \label{eq:opt-adagrad-update}
\end{align}

The symbol $\odot$ denotes element-wise multiplication, and division and
square roots in vector expressions are also element-wise. Coordinates
with large accumulated squared gradients receive smaller effective
learning rates. This behavior is useful when gradient scales are
heterogeneous or sparse. Because the accumulator grows monotonically,
however, the effective learning rates can become very small over long
training runs.

RMSProp replaces the cumulative sum with an exponential moving average:

\begin{equation}
    \mathbf{v}_t
    =
    \beta_2\mathbf{v}_{t-1}
    +(1-\beta_2)(\mathbf{g}_t\odot\mathbf{g}_t),
    \label{eq:opt-rmsprop-state}
\end{equation}

and uses

\begin{equation}
    \boldsymbol{\theta}_{t+1}
    =
    \boldsymbol{\theta}_t
    -
    \eta_t
    \frac{\mathbf{g}_t}{\sqrt{\mathbf{v}_t}+\epsilon}.
    \label{eq:opt-rmsprop-update}
\end{equation}

The exponential average gives greater weight to recent gradient
magnitudes, allowing the coordinate-wise scaling to adapt as the
optimization dynamics change \parencite{Ranganath2026Navigating}.

\paragraph{Adam.}
Adam combines momentum with adaptive second-moment scaling. It maintains
a first-moment estimate and a second-moment estimate:

\begin{align}
    \mathbf{m}_t
    &=
    \beta_1\mathbf{m}_{t-1}
    +(1-\beta_1)\mathbf{g}_t,
    \label{eq:opt-adam-first-moment}\\
    \mathbf{v}_t
    &=
    \beta_2\mathbf{v}_{t-1}
    +(1-\beta_2)(\mathbf{g}_t\odot\mathbf{g}_t).
    \label{eq:opt-adam-second-moment}
\end{align}

Because both estimates are initialized at zero, they are biased toward
zero during early iterations. Adam corrects this initialization bias by
using

\begin{align}
    \widehat{\mathbf{m}}_t
    &=
    \frac{\mathbf{m}_t}{1-\beta_1^t},
    \label{eq:opt-adam-mhat}\\
    \widehat{\mathbf{v}}_t
    &=
    \frac{\mathbf{v}_t}{1-\beta_2^t}.
    \label{eq:opt-adam-vhat}
\end{align}

The parameter update is

\begin{equation}
    \boldsymbol{\theta}_{t+1}
    =
    \boldsymbol{\theta}_t
    -
    \eta_t
    \frac{\widehat{\mathbf{m}}_t}
    {\sqrt{\widehat{\mathbf{v}}_t}+\epsilon}.
    \label{eq:opt-adam-update}
\end{equation}

The first moment provides a smoothed update direction, while the second
moment provides coordinate-wise scaling based on historical gradient
magnitudes. Adam therefore combines two mechanisms that are separate in
ordinary SGD: directional smoothing and adaptive effective learning
rates \parencite{Ranganath2026Navigating}.

\paragraph{AdamW and decoupled weight decay.}
Weight decay regularizes parameter magnitudes. In ordinary gradient
descent, adding an $\ell_2$ penalty
$\frac{\lambda}{2}\|\boldsymbol{\theta}\|_2^2$ to the objective adds
$\lambda\boldsymbol{\theta}$ to its gradient. The resulting update can
be written as

\begin{equation}
    \boldsymbol{\theta}_{t+1}
    =
    (1-\eta_t\lambda)\boldsymbol{\theta}_t
    -
    \eta_t\nabla F(\boldsymbol{\theta}_t).
    \label{eq:opt-l2-decay-sgd}
\end{equation}

For an adaptive optimizer, including the penalty in the gradient causes
the regularization term to be transformed by the adaptive preconditioner.
Its effective strength then depends on the optimizer's coordinate-wise
statistics. AdamW separates the decay operation from the adaptive
gradient update:

\begin{equation}
    \boxed{
    \boldsymbol{\theta}_{t+1}
    =
    (1-\eta_t\lambda)\boldsymbol{\theta}_t
    -
    \eta_t
    \frac{\widehat{\mathbf{m}}_t}
    {\sqrt{\widehat{\mathbf{v}}_t}+\epsilon}.
    }
    \label{eq:opt-adamw-update}
\end{equation}

Here, $\lambda$ is the decoupled weight-decay coefficient. The separation
allows the adaptive update and the regularization strength to be tuned
more independently. AdamW is the principal baseline in the supplied
LLM-optimizer survey and in the later systematic pretraining comparison
\parencite{Ranganath2026Navigating,Wen2026Fantastic}.

The optimizer update is used together with a learning-rate schedule.
A common schedule begins with warmup and then decays the learning rate.
For linear warmup,

\begin{equation}
    \eta_t
    =
    \eta_{\max}\frac{t}{T_{\mathrm{warmup}}},
    \qquad
    0\leq t\leq T_{\mathrm{warmup}}.
    \label{eq:opt-linear-warmup}
\end{equation}

A cosine decay phase can then be written as

\begin{equation}
    \eta_t
    =
    \eta_{\min}
    +
    \frac{\eta_{\max}-\eta_{\min}}{2}
    \left[
        1+
        \cos\left(
            \pi
            \frac{t-T_{\mathrm{warmup}}}
            {T-T_{\mathrm{warmup}}}
        \right)
    \right].
    \label{eq:opt-cosine-decay}
\end{equation}

Warmup is especially relevant to adaptive optimizers because their moment
estimates are initially inaccurate. Large early updates can destabilize
training in large Transformer models. Gradient clipping is another
stability mechanism. For a global clipping threshold $\tau$, the clipped
gradient is

\begin{equation}
    \widetilde{\mathbf{g}}_t
    =
    \mathbf{g}_t
    \min\left(
        1,\frac{\tau}{\|\mathbf{g}_t\|_2}
    \right).
    \label{eq:opt-gradient-clipping}
\end{equation}

When the gradient norm is below $\tau$, the gradient is unchanged; when
it exceeds $\tau$, its norm is reduced to $\tau$. In practice, AdamW
training is also affected by the numerical stabilizer $\epsilon$, the
moment coefficients, parameter-specific weight-decay exclusions, and
the precision used to store parameters and optimizer states
\parencite{Ranganath2026Navigating}.

\subsection{Memory-Efficient and Alternative Update Rules}
\label{subsec:optimization-memory-efficient}

The two moment vectors maintained by AdamW require memory proportional
to the number of model parameters. If each state element is stored in
32-bit floating-point precision, the two vectors require

\begin{equation}
    M_{\mathrm{states}}
    =
    2d\times4
    =
    8d\ \text{bytes}.
    \label{eq:opt-adam-state-memory}
\end{equation}

For a model with one billion parameters, this is approximately
$8$ gigabytes for the first and second moments alone. This excludes
model weights, gradients, activations, temporary buffers, and
distributed-training overhead. A simplified mixed-precision accounting
that includes low-precision parameters and gradients, 32-bit master
parameters, and the two 32-bit Adam states can reach approximately
$16d$ bytes \parencite{Ranganath2026Navigating}.

\paragraph{Adafactor and factored second moments.}
Adafactor reduces the storage required for adaptive second-moment
statistics by exploiting the matrix structure of a parameter tensor.
For a matrix $\mathbf{W}\in\mathbb{R}^{m\times n}$, a full coordinate-wise
second-moment accumulator contains $mn$ entries. Adafactor instead
maintains row and column statistics with $m$ and $n$ entries,
respectively. In a factored approximation, the estimated second-moment
matrix is reconstructed from these lower-dimensional statistics.
The state storage therefore scales as $O(m+n)$ rather than $O(mn)$.
This reduces memory, at the cost of replacing independent statistics
for every matrix entry with a structured approximation
\parencite{Ranganath2026Navigating}.

\paragraph{Low-precision and grouped optimizer states.}
Low-bit optimizer methods reduce memory by representing optimizer
statistics with fewer bits. Block-wise quantization partitions a
state tensor into blocks and quantizes values relative to block-level
scales. The resulting approximation error must be controlled because
the moment estimates directly influence update magnitudes.

Adam-mini reduces the number of independent adaptive statistics by
grouping parameters and sharing adaptive information within groups.
Its motivation is that maintaining an independent learning-rate
statistic for every parameter can be redundant. LOMO takes a different
systems-oriented approach by fusing gradient computation and parameter
updates, reducing the need to retain conventional gradient and optimizer
state tensors simultaneously. These methods reduce different components
of the memory and execution burden, and their trade-offs must be
evaluated in the context of the training implementation
\parencite{Ranganath2026Navigating}.

\paragraph{Lion and sign-based updates.}
Lion replaces adaptive second-moment normalization with a sign-based
update derived from a momentum-like quantity. A representative form is

\begin{align}
    \mathbf{u}_t
    &=
    \beta_1\mathbf{m}_{t-1}
    +(1-\beta_1)\mathbf{g}_t,
    \label{eq:opt-lion-direction}\\
    \boldsymbol{\theta}_{t+1}
    &=
    (1-\eta_t\lambda)\boldsymbol{\theta}_t
    -
    \eta_t\operatorname{sign}(\mathbf{u}_t),
    \label{eq:opt-lion-update}\\
    \mathbf{m}_t
    &=
    \beta_2\mathbf{m}_{t-1}
    +(1-\beta_2)\mathbf{g}_t.
    \label{eq:opt-lion-state}
\end{align}

The sign function is applied element-wise. It returns $-1$, $0$, or $1$
according to the sign of its argument. Unlike AdamW, which rescales
each coordinate according to a second-moment estimate, Lion primarily
retains directional information in the update. It therefore removes
the need for an Adam-style second-moment vector, but its effective
update magnitude is controlled differently and requires its own
learning-rate and weight-decay tuning
\parencite{Ranganath2026Navigating,Wen2026Fantastic}.

\paragraph{Large-batch optimization.}
Large minibatches can improve accelerator utilization and reduce
gradient variance, but the optimization dynamics may change as the
batch size increases. LAMB introduces layer-wise trust ratios to
scale updates according to the relative norms of parameters and
their proposed updates. For a layer with parameters
$\mathbf{w}_{\ell}$ and an Adam-derived update direction
$\mathbf{u}_{\ell}$, the trust ratio has the form

\begin{equation}
    r_{\ell}
    =
    \frac{\|\mathbf{w}_{\ell}\|_2}
         {\|\mathbf{u}_{\ell}\|_2},
    \label{eq:opt-lamb-trust-ratio}
\end{equation}

with practical implementations including stabilizing details for
small denominators and weight decay. The layer-wise scaling changes
the relative magnitude of updates across layers and is intended to
support large-batch training. It also illustrates that an optimizer's
behavior depends on the interaction between its update rule, batch
size, and parameter-group structure
\parencite{Ranganath2026Navigating}.

\subsection{Curvature-Aware and Second-Order Optimization}
\label{subsec:optimization-second-order}

The gradient describes the first-order sensitivity of the objective,
but it does not fully describe the local geometry of the optimization
landscape. For a twice-differentiable objective, the second-order
Taylor approximation is

\begin{equation}
    F(\boldsymbol{\theta}_t+\Delta\boldsymbol{\theta})
    \approx
    F(\boldsymbol{\theta}_t)
    +
    \mathbf{g}_t^{\mathsf{T}}\Delta\boldsymbol{\theta}
    +
    \frac{1}{2}
    \Delta\boldsymbol{\theta}^{\mathsf{T}}
    \mathbf{H}_t
    \Delta\boldsymbol{\theta},
    \label{eq:opt-second-order-taylor}
\end{equation}

where $\mathbf{g}_t=\nabla F(\boldsymbol{\theta}_t)$ and

\begin{equation}
    \mathbf{H}_t
    =
    \nabla_{\boldsymbol{\theta}}^2F(\boldsymbol{\theta}_t)
    \label{eq:opt-hessian}
\end{equation}

is the Hessian. The Hessian is a $d\times d$ matrix whose $(i,j)$ entry
is the second partial derivative
$\partial^2F/(\partial\theta_i\partial\theta_j)$. Its diagonal entries
describe curvature along individual coordinates, while its
off-diagonal entries describe interactions between parameter
coordinates.

If the Hessian is positive definite, minimizing the quadratic
approximation in \cref{eq:opt-second-order-taylor} gives

\begin{equation}
    \mathbf{H}_t\Delta\boldsymbol{\theta}_t
    =
    -\mathbf{g}_t,
    \qquad
    \Delta\boldsymbol{\theta}_t
    =
    -\mathbf{H}_t^{-1}\mathbf{g}_t.
    \label{eq:opt-newton-step}
\end{equation}

This is the Newton step. The inverse Hessian rescales the gradient
according to local curvature. To understand this effect, consider the
eigendecomposition of a symmetric positive-definite Hessian:

\begin{equation}
    \mathbf{H}
    =
    \mathbf{Q}\boldsymbol{\Lambda}\mathbf{Q}^{\mathsf{T}},
    \qquad
    \boldsymbol{\Lambda}
    =
    \operatorname{diag}(\lambda_1,\ldots,\lambda_d).
    \label{eq:opt-hessian-eigendecomposition}
\end{equation}

The columns of $\mathbf{Q}$ are orthonormal curvature directions, and
the eigenvalues describe curvature magnitudes along those directions.
The condition number is

\begin{equation}
    \kappa(\mathbf{H})
    =
    \frac{\lambda_{\max}(\mathbf{H})}
         {\lambda_{\min}(\mathbf{H})}.
    \label{eq:opt-hessian-condition}
\end{equation}

A large condition number indicates that curvature varies substantially
across directions. Gradient descent must use a learning rate small
enough for the steep directions, which can make progress slow in
shallow directions. Newton's method attempts to compensate by scaling
each eigendirection inversely to its curvature.

Exact Newton optimization is infeasible for a model with billions of
parameters. A Hessian with $d^2$ entries requires $O(d^2)$ storage, and
constructing or factorizing it is also computationally prohibitive.
For $d=10^9$, the matrix contains $10^{18}$ entries; at four bytes per
entry, its storage alone would be approximately $4\times10^{18}$ bytes.
Large-scale methods must therefore approximate or avoid explicit
Hessian construction \parencite{Bottou2018Optimization,Ranganath2026Navigating}.

\paragraph{Quasi-Newton and Hessian-free methods.}
Quasi-Newton methods approximate curvature using changes in parameters
and gradients. Define

\begin{align}
    \mathbf{s}_t
    &=
    \boldsymbol{\theta}_{t+1}-\boldsymbol{\theta}_t,
    \label{eq:opt-quasinewton-s}\\
    \mathbf{y}_t
    &=
    \mathbf{g}_{t+1}-\mathbf{g}_t.
    \label{eq:opt-quasinewton-y}
\end{align}

The pair $(\mathbf{s}_t,\mathbf{y}_t)$ contains information about how
the gradient changes along the recent parameter displacement. BFGS
uses such pairs to update an approximation of the inverse Hessian.
Limited-memory BFGS (L-BFGS) retains only a finite history of these
pairs rather than storing a complete dense matrix. Stochastic
minibatches complicate these methods because gradient differences
contain sampling noise, and line searches or synchronization can
introduce additional cost.

Hessian-free methods avoid explicitly forming the Hessian and instead
use iterative procedures to approximate a Newton-like direction.
Bottou, Curtis, and Nocedal review both stochastic quasi-Newton and
Hessian-free approaches as attempts to retain useful curvature
information while controlling computational cost
\parencite{Bottou2018Optimization}.

\paragraph{Curvature estimation in language-model optimizers.}
Sophia is a curvature-aware optimizer discussed in the supplied
LLM-optimizer survey. It combines momentum information with an
estimate of diagonal curvature and clips the resulting
curvature-normalized update. The curvature estimate is intended to
provide information that is not identical to the historical
squared-gradient statistics used by Adam. The clipping operation
limits excessively large coordinate-wise updates. Sophia is an
example of an approach that introduces approximate second-order
information without constructing the complete Hessian
\parencite{Ranganath2026Navigating}.

\subsection{Structured Matrix and Tensor Preconditioning: Shampoo}
\label{subsec:optimization-shampoo}

The previous discussion shows that a full Hessian or full-dimensional
preconditioner is too expensive for large neural networks. Diagonal
adaptive methods avoid this cost by treating coordinates independently,
but this discards correlations between different entries of a parameter
matrix. Transformer layers contain matrix-valued parameters, so a
natural alternative is to retain matrix structure while approximating
the full preconditioner. Shampoo, introduced by Gupta, Koren, and
Singer, follows this approach by maintaining one preconditioning
matrix for each tensor dimension \parencite{Gupta2018Shampoo}.

Consider a matrix parameter
$\mathbf{W}\in\mathbb{R}^{m\times n}$ with gradient
$\mathbf{G}_t\in\mathbb{R}^{m\times n}$. A full preconditioner for the
vectorized gradient would have dimension $mn\times mn$ and require
$O(m^2n^2)$ storage. Shampoo instead maintains a left preconditioner
$\mathbf{L}_t\in\mathbb{R}^{m\times m}$ and a right preconditioner
$\mathbf{R}_t\in\mathbb{R}^{n\times n}$:

\begin{align}
    \mathbf{L}_t
    &=
    \mathbf{L}_{t-1}
    +
    \mathbf{G}_t\mathbf{G}_t^{\mathsf{T}},
    \label{eq:opt-shampoo-left}\\
    \mathbf{R}_t
    &=
    \mathbf{R}_{t-1}
    +
    \mathbf{G}_t^{\mathsf{T}}\mathbf{G}_t.
    \label{eq:opt-shampoo-right}
\end{align}

A positive regularization term is commonly used in the initialization,
for example
$\mathbf{L}_0=\epsilon\mathbf{I}_m$ and
$\mathbf{R}_0=\epsilon\mathbf{I}_n$. The left matrix accumulates
row-wise gradient correlations: its $(i,j)$ entry is the inner product
between rows $i$ and $j$ of the gradient. The right matrix similarly
accumulates column-wise gradient correlations.

Both matrices are positive semidefinite because, for any
$\mathbf{x}\in\mathbb{R}^{m}$,

\begin{equation}
    \mathbf{x}^{\mathsf{T}}
    \mathbf{G}_t\mathbf{G}_t^{\mathsf{T}}\mathbf{x}
    =
    \|\mathbf{G}_t^{\mathsf{T}}\mathbf{x}\|_2^2
    \geq0,
    \label{eq:opt-shampoo-psd}
\end{equation}

and the same argument applies to
$\mathbf{G}_t^{\mathsf{T}}\mathbf{G}_t$. Regularization makes the
preconditioners positive definite, allowing inverse fractional matrix
powers to be defined.

For a symmetric positive-definite matrix
$\mathbf{A}=\mathbf{Q}\boldsymbol{\Lambda}\mathbf{Q}^{\mathsf{T}}$,
a fractional power is defined spectrally:

\begin{equation}
    \mathbf{A}^{\alpha}
    =
    \mathbf{Q}\boldsymbol{\Lambda}^{\alpha}\mathbf{Q}^{\mathsf{T}}.
    \label{eq:opt-matrix-fractional-power}
\end{equation}

In particular,
$\mathbf{A}^{-1/4}
=\mathbf{Q}\operatorname{diag}(\lambda_i^{-1/4})\mathbf{Q}^{\mathsf{T}}$.
This operation rotates into the eigenbasis, rescales each eigendirection,
and rotates back. The original matrix form of Shampoo uses the update

\begin{equation}
    \boxed{
    \mathbf{W}_{t+1}
    =
    \mathbf{W}_t
    -
    \eta_t
    \mathbf{L}_t^{-1/4}
    \mathbf{G}_t
    \mathbf{R}_t^{-1/4}.
    }
    \label{eq:opt-shampoo-matrix-update}
\end{equation}

Left multiplication transforms the row geometry, and right multiplication
transforms the column geometry. Unlike a diagonal optimizer, the
transformation can mix information across coordinates within each matrix
dimension.

The relationship to a full preconditioner can be seen through the
Kronecker product. With a conventional column-major vectorization,
the identity

\begin{equation}
    \operatorname{vec}(\mathbf{A}\mathbf{G}\mathbf{B})
    =
    (\mathbf{B}^{\mathsf{T}}\otimes\mathbf{A})
    \operatorname{vec}(\mathbf{G})
    \label{eq:opt-vec-kronecker}
\end{equation}

implies that the Shampoo matrix update applies the structured
preconditioner

\begin{equation}
    \mathbf{R}_t^{-1/4}\otimes\mathbf{L}_t^{-1/4}
    \label{eq:opt-shampoo-kronecker}
\end{equation}

to the vectorized gradient. The Kronecker product is not explicitly
constructed; its action is obtained by two matrix multiplications.
The storage requirement is therefore $O(m^2+n^2)$ rather than
$O(m^2n^2)$. Exact eigendecompositions or matrix-root computations
nevertheless have costs that scale cubically with the individual
matrix dimensions, which motivates practical approximations and
infrequent updates \parencite{Gupta2018Shampoo,Vyas2025SOAP}.

Shampoo generalizes from matrices to a tensor
$\mathcal{G}_t\in
\mathbb{R}^{n_1\times\cdots\times n_k}$ of order $k$. Let
$\mathbf{G}^{(i)}_t$ denote the matricization of the tensor along
mode $i$, with shape $n_i\times n_{-i}$, where
$n_{-i}=\prod_{j\ne i}n_j$. The preconditioner for mode $i$ is

\begin{equation}
    \mathbf{H}_{t,i}
    =
    \mathbf{H}_{t-1,i}
    +
    \mathbf{G}^{(i)}_t
    \left(\mathbf{G}^{(i)}_t\right)^{\mathsf{T}}.
    \label{eq:opt-shampoo-tensor-preconditioner}
\end{equation}

The gradient is transformed along each tensor mode:

\begin{equation}
    \widetilde{\mathcal{G}}_t
    =
    \mathcal{G}_t
    \times_1\mathbf{H}_{t,1}^{-1/(2k)}
    \times_2\mathbf{H}_{t,2}^{-1/(2k)}
    \cdots
    \times_k\mathbf{H}_{t,k}^{-1/(2k)},
    \label{eq:opt-shampoo-tensor-transform}
\end{equation}

where $\times_i$ denotes multiplication along mode $i$. When $k=2$,
the exponent becomes $-1/4$, recovering the matrix formulation.
The corresponding parameter update is

\begin{equation}
    \mathcal{W}_{t+1}
    =
    \mathcal{W}_t-\eta_t\widetilde{\mathcal{G}}_t.
    \label{eq:opt-shampoo-tensor-update}
\end{equation}

The preconditioner storage is proportional to
$O(\sum_{i=1}^{k}n_i^2)$, rather than the square of the total number
of tensor entries. The method can use diagonal preconditioning for
dimensions whose full matrix preconditioners would be too expensive.
The approximation is structured: it retains correlations within each
tensor dimension but does not represent arbitrary correlations between
different parameter tensors \parencite{Gupta2018Shampoo}.

The original Shampoo analysis is carried out in online convex
optimization. For matrix gradients of rank at most $r$, Gupta, Koren,
and Singer establish a regret bound of the form

\begin{equation}
    \sum_{t=1}^{T} f_t(\mathbf{W}_t)
    -
    \sum_{t=1}^{T} f_t(\mathbf{W}^{*})
    \leq
    \sqrt{2r}\,D\,
    \operatorname{Tr}(\mathbf{L}_T^{1/4})
    \operatorname{Tr}(\mathbf{R}_T^{1/4}),
    \label{eq:opt-shampoo-regret}
\end{equation}

where

\begin{align}
    \mathbf{L}_T
    &=
    \epsilon\mathbf{I}_m+
    \sum_{t=1}^{T}\mathbf{G}_t\mathbf{G}_t^{\mathsf{T}},
    \label{eq:opt-shampoo-LT}\\
    \mathbf{R}_T
    &=
    \epsilon\mathbf{I}_n+
    \sum_{t=1}^{T}\mathbf{G}_t^{\mathsf{T}}\mathbf{G}_t,
    \label{eq:opt-shampoo-RT}\\
    D
    &=
    \max_{t\in\{1,\ldots,T\}}
    \|\mathbf{W}_t-\mathbf{W}^{*}\|_F.
    \label{eq:opt-shampoo-D}
\end{align}

Here, $\mathbf{W}^{*}$ is a comparator matrix and $\|\cdot\|_F$
denotes the Frobenius norm. Under the paper's assumptions, the regret
has the optimal $O(\sqrt{T})$ dependence for online convex
optimization. This is a convex online guarantee; it is not a proof
of global convergence for nonconvex Transformer training
\parencite{Gupta2018Shampoo}.

The original algorithm's matrix-root computations can introduce
substantial overhead if performed at every iteration. Practical
implementations therefore use engineering choices such as periodic
preconditioner-root updates, momentum, and diagonal approximations for
very large dimensions. Shampoo's experiments include a language-model
setting with an attention model of approximately $9.8$ million
parameters on the LM1B benchmark. These experiments demonstrate the
potential of structured preconditioning, but they should not be
interpreted as a direct guarantee of the same speedup at contemporary
billion-parameter scales \parencite{Gupta2018Shampoo}.

\subsection{SOAP: Adam in the Shampoo Eigenbasis}
\label{subsec:optimization-soap}

Shampoo uses matrix-valued preconditioners, but recomputing their
eigenvectors or matrix roots can be expensive. Computing these
quantities less frequently reduces overhead, although the preconditioner
basis then becomes stale between updates. SOAP, introduced by Vyas
and colleagues, addresses this issue by running Adam in the
preconditioner's eigenbasis and continuing to update the Adam
second-moment statistics between basis refreshes
\parencite{Vyas2025SOAP}.

For a matrix-valued layer, let the batch gradient be
$\mathbf{G}_t\in\mathbb{R}^{m\times n}$. Consider the expected
left and right gradient covariance matrices

\begin{align}
    \mathbf{L}
    &=
    \mathbb{E}_{B}
    [\mathbf{G}_{B}\mathbf{G}_{B}^{\mathsf{T}}],
    \label{eq:opt-soap-L}\\
    \mathbf{R}
    &=
    \mathbb{E}_{B}
    [\mathbf{G}_{B}^{\mathsf{T}}\mathbf{G}_{B}],
    \label{eq:opt-soap-R}
\end{align}

where the expectation is over batches. Let their eigendecompositions be

\begin{align}
    \mathbf{L}
    &=
    \mathbf{Q}_{L}\boldsymbol{\Lambda}_{L}
    \mathbf{Q}_{L}^{\mathsf{T}},
    \label{eq:opt-soap-left-eigen}\\
    \mathbf{R}
    &=
    \mathbf{Q}_{R}\boldsymbol{\Lambda}_{R}
    \mathbf{Q}_{R}^{\mathsf{T}}.
    \label{eq:opt-soap-right-eigen}
\end{align}

The rotated gradient is

\begin{equation}
    \mathbf{G}'_t
    =
    \mathbf{Q}_{L}^{\mathsf{T}}
    \mathbf{G}_t
    \mathbf{Q}_{R}.
    \label{eq:opt-soap-rotated-gradient}
\end{equation}

The SOAP paper first establishes a connection between a modified,
idealized Shampoo algorithm using a power of $1/2$ and Adafactor
operating in the eigenbasis supplied by Shampoo. The equivalence
uses dataset-averaged covariance matrices and a scalar correction
to the layer learning rate. This is distinct from the original
Shampoo exponent of $1/4$. The connection motivates replacing
Adafactor with Adam in the rotated coordinate system
\parencite{Vyas2025SOAP}.

In the practical SOAP algorithm, the first-moment state is maintained
in the original parameter coordinates:

\begin{equation}
    \mathbf{M}_t
    =
    \beta_1\mathbf{M}_{t-1}
    +(1-\beta_1)\mathbf{G}_t.
    \label{eq:opt-soap-first-moment}
\end{equation}

It is then rotated into the current preconditioner basis:

\begin{equation}
    \mathbf{M}'_t
    =
    \mathbf{Q}_{L}^{\mathsf{T}}
    \mathbf{M}_t
    \mathbf{Q}_{R}.
    \label{eq:opt-soap-rotated-momentum}
\end{equation}

The second-moment state is maintained in the rotated coordinates:

\begin{equation}
    \mathbf{V}'_t
    =
    \beta_2\mathbf{V}'_{t-1}
    +(1-\beta_2)
    (\mathbf{G}'_t\odot\mathbf{G}'_t).
    \label{eq:opt-soap-second-moment}
\end{equation}

After the appropriate bias corrections, the Adam update in the rotated
space is

\begin{equation}
    \mathbf{N}'_t
    =
    \frac{\widehat{\mathbf{M}}'_t}
    {\sqrt{\widehat{\mathbf{V}}'_t}+\epsilon}.
    \label{eq:opt-soap-adam-rotated}
\end{equation}

The update is transformed back to the original coordinates:

\begin{equation}
    \mathbf{N}_t
    =
    \mathbf{Q}_{L}
    \mathbf{N}'_t
    \mathbf{Q}_{R}^{\mathsf{T}},
    \label{eq:opt-soap-inverse-rotation}
\end{equation}

and the parameters are updated using the AdamW-style decoupled
weight-decay operation:

\begin{equation}
    \mathbf{W}_{t+1}
    =
    (1-\eta_t\lambda)\mathbf{W}_t
    -
    \eta_t\mathbf{N}_t.
    \label{eq:opt-soap-update}
\end{equation}

The left and right covariance estimates are updated online:

\begin{align}
    \mathbf{L}_t
    &=
    \beta_2\mathbf{L}_{t-1}
    +(1-\beta_2)\mathbf{G}_t\mathbf{G}_t^{\mathsf{T}},
    \label{eq:opt-soap-running-L}\\
    \mathbf{R}_t
    &=
    \beta_2\mathbf{R}_{t-1}
    +(1-\beta_2)\mathbf{G}_t^{\mathsf{T}}\mathbf{G}_t.
    \label{eq:opt-soap-running-R}
\end{align}

The eigenvector basis is refreshed every $f$ iterations, where $f$ is
the preconditioning frequency. SOAP uses power iteration followed by
QR decomposition to update eigenvector estimates. In a simplified
power-iteration step, a positive-semidefinite preconditioner
$\mathbf{P}$ and a current basis estimate $\mathbf{Q}$ produce

\begin{equation}
    \mathbf{S}=\mathbf{P}\mathbf{Q},
    \qquad
    \mathbf{Q}_{\mathrm{new}}
    =
    \operatorname{QR}(\mathbf{S}),
    \label{eq:opt-soap-power-qr}
\end{equation}

where $\operatorname{QR}$ denotes the orthogonal factor from a QR
decomposition. The first basis is initialized by an eigendecomposition.
For one-dimensional parameters, SOAP uses ordinary AdamW. For very
large dimensions, a rotation may be fixed to the identity to avoid
prohibitive memory and computation; if both rotations of a matrix
layer are identities, the update reduces to Adam-style optimization
\parencite{Vyas2025SOAP}.

Vyas and colleagues evaluate SOAP on language-model pretraining runs
with approximately $360$-million- and $660$-million-parameter models.
In their large-batch experiments, SOAP reduces the number of
optimization iterations by more than $40\%$ and wall-clock time by
more than $35\%$ relative to AdamW, and improves both measures by
approximately $20\%$ relative to Shampoo. These are results for the
experimental configurations in that paper; they do not establish
that the same improvements hold at all model scales, hardware
configurations, or data-to-model ratios
\parencite{Vyas2025SOAP}.

\subsection{Low-Rank and Matrix-Orthogonalized Optimizers}
\label{subsec:optimization-low-rank-matrix}

Shampoo and SOAP exploit matrix structure through covariance
preconditioning. Other optimizers exploit different structural
properties of neural-network gradients and parameter matrices.
Low-rank methods reduce the dimension in which optimizer states are
maintained, whereas matrix-orthogonalized methods transform a
matrix-valued update according to its singular-value structure
\parencite{Ranganath2026Navigating}.

\paragraph{GaLore and low-rank projection.}
For a gradient matrix
$\mathbf{G}_t\in\mathbb{R}^{m\times n}$, a low-rank approximation
takes the form

\begin{equation}
    \mathbf{G}_t
    \approx
    \mathbf{U}_t\mathbf{S}_t\mathbf{V}_t^{\mathsf{T}},
    \qquad
    r\ll\min(m,n),
    \label{eq:opt-low-rank-svd}
\end{equation}

where $\mathbf{U}_t$ and $\mathbf{V}_t$ contain $r$ basis vectors
and $\mathbf{S}_t$ contains the corresponding reduced-dimensional
coefficients. Projection-based optimizers maintain adaptive states
in a low-dimensional subspace rather than in the full matrix space.
For a right projection matrix
$\mathbf{P}_t\in\mathbb{R}^{n\times r}$, the projected gradient is

\begin{equation}
    \widetilde{\mathbf{G}}_t
    =
    \mathbf{G}_t\mathbf{P}_t.
    \label{eq:opt-low-rank-projection}
\end{equation}

An update computed in the projected space can be mapped back by
multiplication with $\mathbf{P}_t^{\mathsf{T}}$. GaLore uses
gradient-derived low-rank subspaces to reduce the memory required
for adaptive optimizer states. It is important to distinguish this
from parameter-efficient fine-tuning: low-rank optimizer projection
can reduce optimizer-state memory while still updating the original
model parameters, whereas parameter-efficient fine-tuning restricts
the trainable parameterization or subset of parameters
\parencite{Ranganath2026Navigating,Vyas2025SOAP}.

\paragraph{Muon and matrix orthogonalization.}
Muon applies an approximate orthogonalization operation to a
momentum-based matrix update. Let

\begin{equation}
    \mathbf{M}_t
    =
    \beta\mathbf{M}_{t-1}
    +(1-\beta)\mathbf{G}_t
    \label{eq:opt-muon-momentum}
\end{equation}

be a momentum matrix. The update is transformed by an approximate
orthogonalization operator:

\begin{equation}
    \mathbf{U}_t
    =
    \operatorname{NS}(\mathbf{M}_t),
    \qquad
    \mathbf{W}_{t+1}
    =
    \mathbf{W}_t-\eta_t\mathbf{U}_t,
    \label{eq:opt-muon-update}
\end{equation}

where $\operatorname{NS}$ denotes a Newton--Schulz-based
approximation to matrix orthogonalization. To understand the
operation, consider the singular value decomposition

\begin{equation}
    \mathbf{M}
    =
    \mathbf{U}\boldsymbol{\Sigma}\mathbf{V}^{\mathsf{T}}.
    \label{eq:opt-muon-svd}
\end{equation}

The corresponding polar orthogonal factor is
$\mathbf{U}\mathbf{V}^{\mathsf{T}}$. It retains the left and right
singular directions while normalizing the singular values. Direct
singular value decomposition at every update can be expensive;
Newton--Schulz iterations approximate the transformation using
matrix multiplications. Their coefficients and normalization are
chosen so that the singular values move toward the desired range.
Muon therefore differs from Shampoo: Shampoo accumulates covariance
statistics and constructs matrix preconditioners, whereas Muon
orthogonalizes a momentum-derived matrix update
\parencite{Ranganath2026Navigating,Wen2026Fantastic}.

In practical training, matrix-based updates are applied to suitable
dense weight matrices, while embeddings, normalization parameters,
and other parameter groups may use different update rules. This
parameter-group-specific treatment reflects the fact that not every
tensor has the same shape or benefits from the same matrix operation
\parencite{Ranganath2026Navigating}.

\subsection{Systems Constraints and Numerical Implementation}
\label{subsec:optimization-systems-constraints}

The mathematical update equation is only one component of an optimizer's
practical cost. Large language models are commonly trained with mixed
precision and distributed computation, and the optimizer must operate
within the memory and communication limits of the complete training
system.

\paragraph{Optimizer-state memory.}
For an AdamW-style mixed-precision configuration, a simplified
accounting may include low-precision model parameters and gradients,
32-bit master parameters, and two 32-bit moment vectors. This gives
an approximate upper accounting of

\begin{equation}
    M_{\mathrm{training}}
    \approx
    2d+2d+4d+4d+4d
    =
    16d\ \text{bytes}.
    \label{eq:opt-mixed-precision-memory}
\end{equation}

The exact total depends on precision, framework implementation,
parameter sharding, and whether master parameters are stored
separately. The equation excludes activation memory and temporary
buffers. Its purpose is to show why optimizer states become a major
memory component as $d$ grows
\parencite{Ranganath2026Navigating}.

\paragraph{Mixed precision and numerical stability.}
LLM training uses numerical formats such as fp16 and bf16 to reduce
memory and increase throughput. Reduced precision changes the
numerical behavior of optimization. Small gradients may underflow,
large values may overflow, and moment estimates may lose precision.
Loss scaling can mitigate underflow in fp16 training, while bf16's
larger exponent range can reduce the need for explicit scaling.
Optimizer states are often retained in fp32 even when model weights
and activations use lower precision.

The stabilizer $\epsilon$ in Adam-like updates can materially affect
the update when second-moment estimates are small. Gradient clipping
limits rare large gradients, and the precision used for moment
accumulation affects the accuracy of adaptive scaling. Matrix
preconditioners introduce further numerical considerations because
eigenvectors, matrix powers, and covariance estimates are sensitive
to finite-precision arithmetic. Comparisons must therefore specify
precision, clipping, loss scaling, and optimizer-state representation
\parencite{Ranganath2026Navigating}.

\paragraph{Distributed optimization and communication.}
In data-parallel training, each worker computes gradients on a local
microbatch. The gradients must be synchronized, commonly through
all-reduce, before replicas apply the same update. Gradient
synchronization can become a communication bottleneck as model size
and worker count increase. Gradient accumulation, overlapping
communication with backpropagation, bucketed communication,
mixed-precision communication, and optimizer-state sharding are
system-level techniques used to manage this cost.

Optimizer-state sharding distributes auxiliary states such as Adam's
first and second moments across workers. This differs from model
sharding, which partitions parameters and may also partition
gradients. Structured preconditioners and matrix-level optimizers
can introduce additional statistics and synchronization requirements.
Consequently, the cost of an optimizer cannot be inferred solely
from the arithmetic in its update equation. It must be measured
within the distributed layout and hardware configuration in which
the model is trained \parencite{Ranganath2026Navigating}.

\subsection{Empirical Evaluation and Optimizer Selection}
\label{subsec:optimization-evaluation}

An optimizer should not be judged only by the rate at which its
training loss decreases during the early part of a run. Two distinct
efficiency measures are important. Token efficiency describes
validation loss as a function of the number of training tokens:

\begin{equation}
    \mathcal{L}_{\mathrm{val}}
    =
    \mathcal{L}_{\mathrm{val}}(N_{\mathrm{tokens}}),
    \label{eq:opt-token-efficiency}
\end{equation}

whereas wall-clock efficiency describes validation loss as a
function of elapsed training time:

\begin{equation}
    \mathcal{L}_{\mathrm{val}}
    =
    \mathcal{L}_{\mathrm{val}}(T_{\mathrm{wall}}).
    \label{eq:opt-wallclock-efficiency}
\end{equation}

An optimizer can require fewer tokens to reach a target loss while
incurring greater computation per update. Conversely, an optimizer
with lower per-step overhead may need more tokens to reach the same
loss. Memory usage, throughput, and downstream performance are
also relevant. A meaningful comparison therefore considers the
complete training budget rather than one isolated metric
\parencite{Ranganath2026Navigating,Wen2026Fantastic}.

Hyperparameter fairness is equally important. Learning rate,
weight decay, warmup duration, momentum coefficients, clipping
thresholds, and optimizer-specific parameters interact with the
update rule. Applying identical hyperparameters to different
optimizers does not ensure a fair comparison. Wen and colleagues
show that Lion and AdamW, for example, can require substantially
different optimal weight-decay coefficients; their reported
comparison includes an optimum near $0.6$ for Lion and near $0.1$
for AdamW in the studied setting. A poorly tuned AdamW baseline
can therefore exaggerate the apparent speedup of an alternative
optimizer \parencite{Wen2026Fantastic}.

Model scale and data-to-model ratio also affect optimizer rankings.
Wen and colleagues evaluate models from approximately $0.1$ billion
to $1.2$ billion parameters and training regimes corresponding to
$1\times$, $2\times$, $4\times$, and $8\times$ the Chinchilla
data-to-model ratio. Their results show that the speedup of
matrix-based optimizers over a well-tuned AdamW baseline decreases
with model size: the reported advantage falls from approximately
$1.4\times$ for $0.1$-billion-parameter models to approximately
$1.1\times$ for $1.2$-billion-parameter models in the reported
$8\times$ regime. They also find that matrix-based optimizers
consistently outperform scalar-based optimizers in their smaller
model experiments, while the preferred optimizer can change as
the data-to-model ratio increases. Muon performs strongly in
smaller-ratio regimes, whereas SOAP and Kron become more
competitive at larger ratios \parencite{Wen2026Fantastic}.

Intermediate checkpoints can also produce misleading rankings.
As learning rates decay, loss curves for different optimizers can
cross, so an optimizer that appears better early in training may
not be better at the target training budget. The comparison should
therefore be made at the intended end-of-training budget, with
hyperparameters tuned separately for each optimizer and model
setting. These results do not invalidate the improvements reported
in the SOAP experiments; they establish that such improvements
must be interpreted in relation to scale, batch size, training
horizon, data-to-model ratio, hardware, and tuning methodology
\parencite{Vyas2025SOAP,Wen2026Fantastic}.

\begin{figure}[htbp]
    \centering
    % ============================================================================
% figures/optimizer_taxonomy.tex
% Included from the section body inside a figure environment.
% Uses standard LaTeX boxes and tabular environments; no TikZ dependency.
% ============================================================================
\begin{center}
\renewcommand{\arraystretch}{1.18}
\setlength{\tabcolsep}{3pt}
\begin{tabular}{c}
\fbox{\parbox{0.78\linewidth}{
    \centering
    \textbf{Stochastic gradient information}\\
    Current minibatch gradient and, where available, historical statistics
}}\\[0.45em]
$\Downarrow$\\[0.45em]
\begin{tabular}{cccc}
\fbox{\parbox{0.215\linewidth}{
    \centering
    \textbf{Scalar / diagonal}\\[0.25em]
    SGD\\
    Momentum\\
    AdaGrad\\
    RMSProp\\
    AdamW\\
    Lion
}}
&
\fbox{\parbox{0.215\linewidth}{
    \centering
    \textbf{Matrix / tensor}\\[0.25em]
    Shampoo\\
    SOAP\\
    Kron\\
    Muon
}}
&
\fbox{\parbox{0.215\linewidth}{
    \centering
    \textbf{Memory-oriented}\\[0.25em]
    Adafactor\\
    Adam-mini\\
    Low-bit states\\
    LOMO\\
    GaLore
}}
&
\fbox{\parbox{0.215\linewidth}{
    \centering
    \textbf{Curvature / projection}\\[0.25em]
    Sophia\\
    Quasi-Newton\\
    Hessian-free\\
    Low-rank projection
}}
\end{tabular}\\[0.45em]
$\Downarrow$\\[0.45em]
\fbox{\parbox{0.78\linewidth}{
    \centering
    \textbf{Parameter update}\\
    A transformed gradient-derived direction is applied under a learning-rate
    schedule and the chosen regularization and numerical-stability rules.
}}
\end{tabular}
\end{center}

    \caption{A conceptual taxonomy of optimizer families discussed in
    this section. The categories distinguish coordinate-wise adaptive
    updates, structured matrix preconditioning, memory-oriented
    methods, and curvature or projection methods. They are not
    mutually exclusive: an optimizer can combine more than one
    design principle.}
    \label{fig:optimizer-taxonomy}
\end{figure}
% Alt text for accessibility register: A gradient input branches into
% four groups of optimizer methods: scalar or diagonal methods, matrix
% and tensor preconditioners, memory-efficient methods, and curvature
% or projection-based methods.

\begin{figure}[htbp]
    \centering
    % ============================================================================
% figures/shampoo_soap_geometry.tex
% Included from the section body inside a figure environment.
% Uses standard LaTeX boxes and tabular environments; no TikZ dependency.
% ============================================================================
\begin{center}
\renewcommand{\arraystretch}{1.2}
\setlength{\tabcolsep}{5pt}
\begin{tabular}{c}
\fbox{\parbox{0.74\linewidth}{
    \centering
    \textbf{Matrix gradient}
    $\mathbf{G}_t\in\mathbb{R}^{m\times n}$
}}\\[0.45em]
$\swarrow\hspace{3.2cm}\searrow$\\[0.35em]
\begin{tabular}{ccc}
\fbox{\parbox{0.30\linewidth}{
    \centering
    \textbf{Left covariance}\\
    $\mathbf{L}_t
    \leftarrow
    \mathbf{L}_{t-1}+\mathbf{G}_t\mathbf{G}_t^{\mathsf T}$
}}
&
\hspace{0.8cm}
&
\fbox{\parbox{0.30\linewidth}{
    \centering
    \textbf{Right covariance}\\
    $\mathbf{R}_t
    \leftarrow
    \mathbf{R}_{t-1}+\mathbf{G}_t^{\mathsf T}\mathbf{G}_t$
}}
\end{tabular}\\[0.45em]
$\Downarrow$\\[0.35em]
\begin{tabular}{cc}
\fbox{\parbox{0.40\linewidth}{
    \centering
    \textbf{Shampoo}\\[0.2em]
    $\mathbf{L}_t^{-1/4}
    \mathbf{G}_t
    \mathbf{R}_t^{-1/4}$\\[0.2em]
    Direct left--right matrix preconditioning
}}
&
\fbox{\parbox{0.40\linewidth}{
    \centering
    \textbf{SOAP}\\[0.2em]
    $\mathbf{G}'_t
    =
    \mathbf{Q}_L^{\mathsf T}
    \mathbf{G}_t
    \mathbf{Q}_R$\\[0.2em]
    Adam-style update in the rotated basis,
    followed by inverse rotation
}}
\end{tabular}
\end{tabular}
\end{center}

    \caption{Conceptual comparison of Shampoo and SOAP for a matrix-valued
    parameter. Shampoo applies left and right covariance-derived matrix
    preconditioners. SOAP rotates the gradient into their eigenbasis,
    applies Adam-style adaptive updates in that basis, and rotates the
    update back to the original coordinates.}
    \label{fig:shampoo-soap-geometry}
\end{figure}
% Alt text for accessibility register: The matrix gradient is used to
% construct left and right covariance matrices. Shampoo applies inverse
% fractional powers of those matrices to the gradient. SOAP instead
% rotates the gradient into the covariance eigenbasis, applies Adam
% there, and rotates the update back.

The relationships summarized in \cref{fig:optimizer-taxonomy} and
\cref{fig:shampoo-soap-geometry} show that optimizer design is a
choice about which information from the gradient history and
parameter geometry is retained, and how that information is used
to construct an update. The empirical results show why no single
optimizer can be declared universally superior from one experiment.
The appropriate choice depends on the model scale, the data budget,
the hardware and distributed-training layout, the memory budget,
and the quality of hyperparameter tuning
\parencite{Ranganath2026Navigating,Wen2026Fantastic}.

\subsection{Limitations and Connection to the Next Topic}
\label{subsec:optimization-limitations-transition}

The progression from stochastic gradient descent to adaptive,
memory-efficient, curvature-aware, and matrix-based optimizers
reflects a progression in the information used to construct
parameter updates. SGD uses the current stochastic gradient.
Momentum retains a smoothed history of gradients. AdamW combines
first- and second-moment estimates with decoupled weight decay.
Memory-efficient methods reduce the cost of storing optimizer
statistics, while curvature-aware and structured methods attempt
to represent additional information about the optimization
landscape.

Shampoo demonstrates how matrix and tensor structure can be
retained without explicitly constructing a full-dimensional
preconditioner. SOAP combines this structured coordinate system
with Adam-style adaptive statistics. GaLore uses low-rank
subspaces to reduce optimizer-state memory, and Muon applies
matrix orthogonalization to momentum-derived updates. These
methods address different limitations and should be compared
using both mathematical and system-level criteria
\parencite{Gupta2018Shampoo,Vyas2025SOAP,Ranganath2026Navigating}.

The theoretical guarantees available for these algorithms also
have different scopes. Stochastic-gradient convergence results
depend on assumptions such as smoothness, convexity or
nonconvex stationarity conditions, and bounded stochastic
variance. Shampoo's regret guarantee is established in a
stochastic convex online setting. Empirical improvements on
nonconvex language-model pretraining tasks are important, but
they do not replace these assumptions with a general proof of
global optimality. Similarly, empirical speedups measured at
one scale or data budget do not establish the same speedup at
another scale \parencite{Bottou2018Optimization,Gupta2018Shampoo,
Vyas2025SOAP,Wen2026Fantastic}.

The central conclusion is that an optimizer is not simply a
rule for subtracting a scaled gradient. It is a numerical
procedure for deciding which statistical and geometric
properties of the gradient should influence the next parameter
update, subject to the memory, precision, computation, and
communication constraints of the training system. Having
established how the parameters are updated, the next discussion
can turn to the objective functions that define what the
optimization process is attempting to achieve. The distinction
between the objective and the algorithm that minimizes it is
necessary for understanding the complete training process of
language models.
